
\documentclass[12pt]{amsart}

%==================== Packages ====================
\usepackage[margin=1in]{geometry}
\usepackage{amsmath,amssymb,amsthm}
\usepackage{mathtools}              % paired delimiters, \coloneqq, etc.
\usepackage{enumitem}               % control list labels
\usepackage{bm}                     % bold math: \bm{v}
\usepackage{xcolor}
\usepackage[colorlinks=true,linkcolor=blue,urlcolor=blue]{hyperref}  % load last

\allowdisplaybreaks                 % let long align blocks break across pages
\setlist[enumerate,1]{label=(\alph*)}  % problem parts as (a), (b), (c)

%==================== Number systems ====================
\newcommand{\N}{\mathbb{N}}
\newcommand{\Z}{\mathbb{Z}}
\newcommand{\Q}{\mathbb{Q}}
\newcommand{\R}{\mathbb{R}}
\newcommand{\C}{\mathbb{C}}
\newcommand{\F}{\mathbb{F}}

%==================== Linear algebra ====================
\DeclareMathOperator{\Ran}{Ran}         % range (Treil); \ker already exists
\DeclareMathOperator{\rank}{rank}
\DeclareMathOperator{\nullity}{nullity}
\DeclareMathOperator{\spn}{span}        % NOT \span: that's a TeX primitive
\DeclareMathOperator{\tr}{tr}
\DeclareMathOperator{\col}{Col}
\DeclareMathOperator{\row}{Row}
\DeclareMathOperator{\proj}{proj}
\DeclareMathOperator{\id}{id}
\newcommand{\Mat}[2]{M_{#1\times #2}(\F)}   % \Mat{m}{n} -> M_{m x n}(F)
\newcommand{\Pn}[1][n]{\mathcal{P}_{#1}}    % \Pn -> P_n, \Pn[3] -> P_3
\newcommand{\tp}{^{\mathsf{T}}}             % A\tp -> A^T
\newcommand{\inv}{^{-1}}                    % A\inv -> A^{-1}
\newcommand{\mat}[1]{\begin{bmatrix}#1\end{bmatrix}}
\newcommand{\pmat}[1]{\begin{pmatrix}#1\end{pmatrix}}

%==================== Analysis / topology ====================
\DeclareMathOperator{\diam}{diam}
\DeclareMathOperator{\dist}{dist}
\DeclareMathOperator{\interior}{int}    % NOT \int: that's the integral sign
\DeclareMathOperator{\sgn}{sgn}
\DeclareMathOperator{\supp}{supp}
\newcommand{\closure}[1]{\overline{#1}}
\newcommand{\pre}{^{\mathrm{pre}}}          % f\pre(V), Pugh's preimage
\newcommand{\nbhd}[2]{N_{#1}(#2)}           % \nbhd{r}{p} -> N_r(p), Rudin
\newcommand{\limn}{\lim_{n\to\infty}}
\newcommand{\eps}{\varepsilon}
\newcommand{\dd}{\,\mathrm{d}}              % \int f(x)\dd x
\newcommand{\restr}[2]{\left.#1\right|_{#2}}  % \restr{f}{A} -> f|_A
\newcommand{\seq}[2][n]{\left(#2\right)_{#1\in\N}}  % \seq{p_n} -> (p_n)_{n in N}

%==================== Delimiters (mathtools) ====================
% \abs{x} normal size; \abs*{x} auto-sizes; \abs[\big]{x} manual size
\DeclarePairedDelimiter{\abs}{\lvert}{\rvert}
\DeclarePairedDelimiter{\norm}{\lVert}{\rVert}
\DeclarePairedDelimiter{\floor}{\lfloor}{\rfloor}
\DeclarePairedDelimiter{\ceil}{\lceil}{\rceil}
\DeclarePairedDelimiter{\inner}{\langle}{\rangle}
\DeclarePairedDelimiterX{\set}[1]{\lbrace}{\rbrace}{#1}
\DeclarePairedDelimiterX{\setof}[2]{\lbrace}{\rbrace}{#1 \,:\, #2}  % \setof{x\in\R}{x>0}

%==================== Logic shortcuts ====================
\newcommand{\contra}{\Rightarrow\!\Leftarrow}   % contradiction marker
\newcommand{\st}{\text{ such that }}

%==================== Theorem environments ====================
\theoremstyle{plain}
\newtheorem*{theorem*}{Theorem}
\newtheorem*{lemma*}{Lemma}
\newtheorem*{claim*}{Claim}
\theoremstyle{definition}
\newtheorem*{definition*}{Definition}
\theoremstyle{remark}
\newtheorem*{remark*}{Remark}

\renewcommand{\qedsymbol}{$\blacksquare$}       % black box

% Problem header using the textbook's exercise number: \begin{problem}{2.5}
\newenvironment{problem}[1]{\par\medskip\noindent\textbf{#1.}\ }{\par\smallskip}
% For computational answers that aren't proofs
\newenvironment{solution}{\begin{proof}[Solution]}{\end{proof}}

%==================== Document ====================
\title{Homework 6}
\author{Antony Perez}
\date{\today}

\begin{document}
\maketitle
\begin{problem}{8.1}
    True or false:
    \begin{enumerate}
        \item[a)] Every change of coordinate matrix is square;
        \item[b)] Every chance of coordinate matrix is invertible;
        \item[c)] The matrices $A$ and $B$ are called similar if $B = Q^T AQ$ for some matrix $Q$;
        \item[d)] The matrices $A$ and $B$ are called similar if $B = Q^{-1}AQ$ for some matrix $Q$;
        \item[e)] Similar matrices do not need to be square.    
    \end{enumerate}
\end{problem}

\begin{solution}
    TTFTF
\end{solution}

\begin{problem}{8.2}
    Consider the system of vectors
    \[(1,2,1,1)^T, (0,1,3,1)^T, (0,3,2,0)^T, (0,1,0,0)^T.\]
    \begin{enumerate}
        \item[a)] Prove that it is a basis in $\F^4$. Try to do minimal amount of computations.
        \item[b)] Find the change of coordinate matrix that changes the coordinates in this basis to the standard coordinates in $\F^4$ (i.e to the coordinates in the standard basis $\mathbf{e}_1,\ldots,\mathbf{e}_4$). 
    \end{enumerate}
\end{problem}

\begin{proof}
    Let $v_1,v_2,v_3,v_4$ be the vectors shown above respectively. Since they form a $4\times 4$ matrix, it is sufficient to show that $av_1+bv_2+cv_3+dv_4 = 0$ implies linear independence. We see $a=0$ for the first row; $a+b = 0 \implies b = 0$ for the 4th row; $a + 3b + 2c = 0 \implies c =0$ for row 3; $2a + b + 3c + d =0 \implies c = 0$ for 2nd row. Thus, these vectors are a basis for $\F^4$. 
    \[[I]_{\mathcal{S}\mathcal{B}} = \begin{pmatrix}
        1 & 0 & 0 & 0 \\ 
        2 & 1 & 3 & 1 \\ 
        1 & 3 & 2 & 0 \\
        1 & 1 & 0 & 0 \\ 
    \end{pmatrix}\]
\end{proof}

\begin{problem}{8.3}
    Find the change of coordinates matrix that changes the coordinates in the basis $1,1+t$ in $\mathbb{P}_1$ to the coordinates in the basis $1-t,2t$.
\end{problem}

\begin{solution}
Let $\mathcal{A} = \{1,1+t\}$ and $\mathcal{B} = \{1-t,2t\}$. $1 = a(1-t) + b(2t) \implies a = 1, b=1/2.$ $1+t = a(1-t)+b(2t) \implies a=1,b =1$. Hence,
\[[I]_{\mathcal{B}\mathcal{A}} = \begin{pmatrix}
    1 & 1 \\ \frac{1}{2} & 1
\end{pmatrix}\]
\end{solution}

\begin{problem}{8.4}
    Let $T$ be the linear operator in $F^2$ defined by 
    \[T \begin{pmatrix}
        x \\ y
    \end{pmatrix} = \begin{pmatrix}
        3x+y \\ x-2y
    \end{pmatrix}\]
    Find the matrix of $T$ in the standard basis and in the basis \[(1,1)^T, (1,2)^T\]
\end{problem}

\begin{solution}
    Let $[T]_\mathcal{S} = \begin{pmatrix} 3 & 1 \\ 1 & -2   \end{pmatrix}$. Let $\mathcal{B} = \{(1,1)^T,(1,2)^T\}$. Let $\mathcal{S} = [I]_{\mathcal{S}\mathcal{B}} = \begin{pmatrix}
        1 & 1 \\ 1 & 2
    \end{pmatrix} \implies \mathcal{S}^{-1} = \begin{pmatrix}
        2 & -1 \\ -1 & 1
    \end{pmatrix}$. Thus $[T]_\mathcal{B} = \mathcal{S}^{-1}[T]_\mathcal{S} \mathcal{S} = \begin{pmatrix}
        9 & 13 \\ -5 & -8
    \end{pmatrix}$
\end{solution}

\begin{problem}{8.5}
    Prove, that if $A$ and $B$ are similar matrices, then $\tr A = \tr B$. 
\end{problem}

\begin{proof}
    $B = Q^{-1}AQ$ by similarity. $\tr(XY) = \tr(YX)$ if both products are defined. Let $X = Q^{-1}A$ and $Y = Q$: $\tr(B)=\tr((Q^{-1}A)Q) = \tr(Q(Q^{-1}A)) = \tr(A)$. Thus, $\tr (A) = \tr (B)$. 
\end{proof}

\begin{problem}{8.6}
    Are the matrices 
    \[\begin{pmatrix}
        1 & 3 \\ 2 & 2
    \end{pmatrix} \text{ and } \begin{pmatrix}
        0 & 2 \\ 4 & 2
    \end{pmatrix}\]
    similar? Justify.
\end{problem}

\begin{proof}
    No. Let $A = \begin{pmatrix}
        1 & 3 \\ 2 & 2
    \end{pmatrix}$ and $B = \begin{pmatrix}
        0 & 2 \\ 4 & 2
    \end{pmatrix}$. If $B = Q^{-1}AQ$ then $AQ = QB$. Letting $Q = \begin{pmatrix}
        a & b \\ c & d
    \end{pmatrix},$ we can compare entries and see that $AQ$ and $QB$ force $a=b=c=d$, so $Q$ is not invertible due to $\det Q = 0$. Hence, the matrices aren't similar. 
\end{proof}

\end{document}
